\chapter{Managing Data Frames}
\label{chap:managing-data-frames}

\chapterguide
  {You should be able to create, inspect, and index data frames.}
  {import, inspect, clean, rename, filter, sort, and export tabular data.}

\section{Packages used in Lab 7a}

Lab 7 uses three packages:

\begin{dstable}
\begin{tabularx}{\linewidth}{@{}p{3.0cm}X@{}}
\toprule
\rowcolor{DSPaleBlue!92}
Package & Purpose \\
\midrule
\texttt{readxl} & Import the Excel file \texttt{uforeport.xls}. \\
\texttt{dplyr} & Filter, pipe, and arrange data frames. \\
\texttt{xlsx} & Export an Excel file. \\
\bottomrule
\end{tabularx}
\end{dstable}

Install packages once; load them in each session:

\begin{lstlisting}[style=dsr]
install.packages("readxl")
library(readxl)

any(grepl("readxl", installed.packages()))
\end{lstlisting}

The final line checks whether \texttt{readxl} appears among installed packages.

\section{Working directory and file placement}

Keep source files in the RStudio project and confirm the active directory:

\begin{lstlisting}[style=dsr]
getwd()
list.files()
\end{lstlisting}

Place \texttt{uforeport.xls} and \texttt{titanic.csv} in that project directory.

\begin{keyidea}
If import fails, check \texttt{getwd()} and \texttt{list.files()} before changing the analysis.
\end{keyidea}

\section{Importing Excel and CSV files}

Import Excel with \texttt{read\_excel()}:

\begin{lstlisting}[style=dsr]
uforeport <- read_excel("uforeport.xls")
View(uforeport)
\end{lstlisting}

Import CSV with base R:

\begin{lstlisting}[style=dsr]
titanic <- read.csv("titanic.csv")
View(titanic)
\end{lstlisting}

\section{Identifying missing values}

\texttt{is.na()} identifies missing values.

\subsection{Count all missing values}

\begin{lstlisting}[style=dsr]
sum(is.na(uforeport))
sum(is.na(titanic))
\end{lstlisting}

\subsection{Find their positions}

\begin{lstlisting}[style=dsr]
which(is.na(uforeport))
which(is.na(titanic))
\end{lstlisting}

A small vector illustrates both operations:

\begin{lstlisting}[style=dsr]
demo <- c(1, 2, NA, 4, NA, 6, 7)

sum(is.na(demo))
which(is.na(demo))
\end{lstlisting}

\subsection{Count missing values by column}

\begin{lstlisting}[style=dsr]
print(sapply(uforeport, function(x) sum(is.na(x))))
print(sapply(titanic, function(x) sum(is.na(x))))
\end{lstlisting}

This returns one missing-value count per column.

\section{Removing rows with missing values}

Compare dimensions before and after removing incomplete rows:

\begin{lstlisting}[style=dsr]
dim(uforeport)
uforeport_cleaned <- na.omit(uforeport)
dim(uforeport_cleaned)

dim(titanic)
titanic_cleaned <- na.omit(titanic)
dim(titanic_cleaned)
\end{lstlisting}

The difference shows how many rows were removed.

\begin{pitfall}
\texttt{na.omit()} removes every row containing at least one missing value. Confirm that this loss is acceptable before using the cleaned data for analysis.
\end{pitfall}

\section{Column names and renaming}

\texttt{names()} and \texttt{colnames()} both retrieve data-frame column names:

\begin{lstlisting}[style=dsr]
colnames(uforeport_cleaned)
colnames(titanic_cleaned)

names(uforeport_cleaned)
names(titanic_cleaned)
\end{lstlisting}

Replace spaces with underscores before filtering:

\begin{lstlisting}[style=dsr]
names(uforeport_cleaned) <- gsub(
  " ", "_", colnames(uforeport_cleaned)
)
\end{lstlisting}

For example, \texttt{Colors Reported} becomes \texttt{Colors\_Reported}.

\section{Filtering data}

Use \texttt{dplyr::filter()} for exact and numeric conditions:

\begin{lstlisting}[style=dsr]
install.packages("dplyr")
library(dplyr)

print(filter(uforeport_cleaned, City == "Belton"))
print(filter(uforeport_cleaned, Colors_Reported == "GREEN"))
print(filter(titanic_cleaned, sex == "female"))
print(filter(titanic_cleaned, fare > 500))
\end{lstlisting}

\subsection{Multiple conditions}

Multiple conditions may be written explicitly or separated inside \texttt{filter()}:

\begin{lstlisting}[style=dsr]
print(filter(
  titanic_cleaned,
  sex == "female" & fare > 500
))

# Alternative dplyr form
titanic_cleaned %>%
  filter(sex == "female", fare > 500)

ufo_green <- uforeport_cleaned %>%
  filter(Colors_Reported == "GREEN")
\end{lstlisting}

\section{Arranging rows}

Sort rows with \texttt{arrange()}:

\begin{lstlisting}[style=dsr]
# Ascending
titanic_sortbyfare <- arrange(titanic_cleaned, fare)

# Descending
titanic_sortbyfare <- arrange(titanic_cleaned, desc(fare))
\end{lstlisting}

Because the variable is reassigned, only the descending result remains after both lines.

\section{Exporting analysis results}

Export the transformed results:

\begin{lstlisting}[style=dsr]
install.packages("xlsx")
library(xlsx)
write.xlsx(ufo_green, "ufo_green.xlsx")

write.csv(
  titanic_sortbyfare,
  "titanic_sortbyfare.csv"
)
\end{lstlisting}

The complete workflow is \textbf{import $\rightarrow$ inspect $\rightarrow$ clean $\rightarrow$ transform $\rightarrow$ export}.

\begin{quickcheck}
\begin{enumerate}
  \item Which function tells you the current working directory?
  \item Which import function is used for \texttt{uforeport.xls}, and which for \texttt{titanic.csv}?
  \item How do \texttt{sum(is.na(x))} and \texttt{which(is.na(x))} answer different questions?
  \item What evidence does the handout collect before and after \texttt{na.omit()}?
  \item Which \texttt{dplyr} function sorts rows, and how is descending order requested?
\end{enumerate}
\end{quickcheck}

\begin{chapterrecap}
\begin{itemize}
  \item Project-relative files make imports reproducible.
  \item Inspect missingness before removing incomplete rows.
  \item Rename, filter, sort, and export only after confirming the data structure.
\end{itemize}
\end{chapterrecap}
